% =============================================================================
% 7. ILLUSTRATIVE NUMERICAL EXAMPLE
% Target: 1.5 pages (~750 words)
% Source: appendix_theory.tex / cap4 sec:s3-example (numerical example)
% Status: SKELETON
%
% Purpose: NOT empirical results in real repositories.
% This is a closed-form proof-of-concept that the analytical derivation
% and the Monte Carlo simulation produce the same answer on a controlled
% 6-state SDLC chain. Validates Property prop:convergence-rate and
% Theorem thm:consistency on a verifiable instance.
% =============================================================================

\section{Illustrative Numerical Example}
\label{sec:numerical-example}

\textbf{[TODO — 1.5 p.]}
This section instantiates PATHCAST on a controlled six-state SDLC chain
to demonstrate (i) the closed-form analytical derivation of $N$, $B$
and the expected lead time and (ii) agreement with the Monte Carlo
simulation of Stage~4. The example uses synthetic parameters chosen to
mirror typical SDLC patterns; no real-repository data is involved.

% ---------- §7.1 Setup ----------
\subsection{Setup}
\label{sec:num-setup}

\textbf{[TODO — describe the 6-state chain.]}
States: $S_T = \{\texttt{backlog}, \texttt{in\_progress},
\texttt{review}, \texttt{testing}\}$,
$S_A = \{\texttt{done}, \texttt{cancelled}\}$.

Transition matrix $P$ (with rework edges from \texttt{review} and
\texttt{testing} back to \texttt{in\_progress}):
\[
P =
\begin{pmatrix}
0.0 & 0.9 & 0.0 & 0.0 & 0.0 & 0.1 \\
0.0 & 0.0 & 0.8 & 0.0 & 0.0 & 0.2 \\
0.0 & 0.3 & 0.0 & 0.6 & 0.1 & 0.0 \\
0.0 & 0.2 & 0.0 & 0.0 & 0.8 & 0.0 \\
0.0 & 0.0 & 0.0 & 0.0 & 1.0 & 0.0 \\
0.0 & 0.0 & 0.0 & 0.0 & 0.0 & 1.0 \\
\end{pmatrix}
\]

Per-state mean sojourn $\bar{d}_s$ (in days):
\textbf{[TODO — fill values, e.g.\ backlog 3, in\_progress 4,
review 2, testing 3.]}

% ---------- §7.2 Closed-form analytical computation ----------
\subsection{Closed-Form Computation}
\label{sec:num-analytical}

\textbf{[TODO — compute $Q, R, N=(I-Q)^{-1}, B=NR$ explicitly with
the matrices above. Show one or two intermediate matrices, then the
final $N$ and $B$.]}

Expected lead time from $s_0 = \texttt{backlog}$:
\begin{equation}
\label{eq:num-lead}
\mathbb{E}[\text{LeadTime} \mid s_0 = \texttt{backlog}]
  = \sum_{j \in S_T} N_{\texttt{backlog},j} \cdot \bar{d}_j
  = \textbf{[TODO — numerical value, e.g.\ 14.3 days.]}
\end{equation}

Completion probability:
\begin{equation}
\label{eq:num-completion}
B_{\texttt{backlog},\texttt{done}}
  = \textbf{[TODO — numerical value, e.g.\ 0.84.]}
\end{equation}

% ---------- §7.3 Monte Carlo cross-validation ----------
\subsection{Monte Carlo Cross-Validation}
\label{sec:num-mc}

\textbf{[TODO — simulate $M_{\text{sim}} = 10{,}000$ trajectories of
$P$ with exponential sojourns matched to $\bar{d}_s$. Report:
$\hat{Q}_{0.5}$, $\hat{Q}_{0.85}$, $\hat{P}(\text{Done})$ and the
empirical mean. Compare with closed-form values.]}

Table~\ref{tab:num-comparison} compares analytical and Monte Carlo
estimates. The relative error is below 1\% on every quantity, as
predicted by Property~\ref{prop:convergence-rate}.

\begin{table}[ht]
\centering
\caption{Analytical vs.\ Monte Carlo estimates on the six-state SDLC
chain. $M_{\text{sim}} = 10{,}000$.}
\label{tab:num-comparison}
\small
\begin{tabular}{@{}lrrl@{}}
\toprule
\textbf{Quantity} & \textbf{Analytical} &
\textbf{Monte Carlo} & \textbf{Rel.\ error} \\
\midrule
$\mathbb{E}[\text{LeadTime}]$  & \textbf{[TODO]} & \textbf{[TODO]} & $< 1\%$ \\
$\hat{Q}_{0.5}$                & \textbf{[TODO]} & \textbf{[TODO]} & $< 1\%$ \\
$\hat{P}(\text{Done})$         & \textbf{[TODO]} & \textbf{[TODO]} & $< 1\%$ \\
\bottomrule
\end{tabular}
\end{table}

% ---------- §7.4 What the example illustrates ----------
\subsection{Discussion}
\label{sec:num-discussion}

\textbf{[TODO — 1 paragraph.]}
The example illustrates three properties of the pipeline.
\textbf{(D1)} The fundamental matrix $N$ exposes the contribution of
each transient state to the expected lead time, locating
\texttt{review} as the dominant contributor through the rework loop.
\textbf{(D2)} The quasi-stationary distribution
(Equation~\textbf{[TODO ref]}) identifies \texttt{review} as the
accumulating state, a bottleneck diagnosis unavailable to
throughput-based Monte Carlo.
\textbf{(D3)} Analytical and Monte Carlo answers agree within 1\%, a
controlled cross-validation that the simulation implements the same
model as the closed-form derivation.

\textbf{This example is not an empirical evaluation.} The empirical
evaluation across 48 open-source repositories is the subject of the
companion paper (Section~\ref{sec:evaluation-plan}).
